%%
% The BIThesis Template for Bachelor Graduation Thesis
%
% 北京理工大学毕业设计（论文）致谢 —— 使用 XeLaTeX 编译
%
% Copyright 2020-2023 BITNP
%
% This work may be distributed and/or modified under the
% conditions of the LaTeX Project Public License, either version 1.3
% of this license or (at your option) any later version.
% The latest version of this license is in
%   https://www.latex-project.org/lppl.txt
% and version 1.3 or later is part of all distributions of LaTeX
% version 2005/12/01 or later.
%
% This work has the LPPL maintenance status `maintained'.
%
% The Current Maintainer of this work is Feng Kaiyu.
%
% Compile with: xelatex -> biber -> xelatex -> xelatex

% 致谢部分尽量不使用 \subsection 二级标题，只使用 \section 一级标题
\begin{acknowledgements}
行文至此，落笔为终。四载春秋，转瞬成诗。
那个初入校园、满怀憧憬与热望的青涩少年，如今已伫立在毕业的渡口。倏忽四载，是导师的教诲与同窗的笑语交汇成的生命诗篇，亦是屏幕前日升月落的坚守与图书馆的星辉点点织就成的青春画卷。纵有课题攻关的困顿与论文改稿的辗转焦灼，但每一份挫折都在你们的陪伴下化作破茧成蝶的勇气。愿此去经年，我们都能在各自的山海间，续写这场始于热爱的奔赴。

桃李不言，下自成蹊。

首先，特别感谢我的毕设导师陆慧梅老师。在整个毕业设计过程中，陆老师始终给予我极大的支持和悉心的指导。无论是开题答辩时的方向修正，中期检查时的进度把控，还是论文终稿的字斟句酌，陆老师都倾注了大量心血；从选题的确定、方向的把控，再到细节的打磨，陆老师始终以极大的耐心给予专业细致的指导。陆老师严谨治学、踏实认真的态度深深感染了我，为我的毕业论文奠定了坚实基础。

同时，我也要特别感谢向勇老师。在每周例行的组会中，向老师都会耐心听取我的汇报，针对我所遇到的问题提出建设性的意见和改进方案。虽然我眼中的向老师较为严格，但正是这种指导方式不断鞭策我持续向前。在一次次讨论和思维碰撞中，向老师敏锐的洞察力和深厚的专业素养使我在毕设过程中少走了许多弯路，对此我将永远铭记在心。

在毕设过程中，我非常幸运地得到了许多优秀学长的帮助。

我要特别感谢廖东海学长，他在我选题之初就给予了极大的帮助，并在整个研究过程中一路指导我。每当我遇到难题时，廖东海学长总是耐心的倾听解答，不厌其烦地为我指点迷津。他的认真态度和严谨精神是我学习的榜样。

同时，我也要感谢赵方亮学长。他为我提供了关于TAIC方面的技术支持和宝贵的硬件资源，并在我遇到研究难题时耐心指导、细心解答，帮助我理清思路、突破瓶颈。他的无私帮助和专业精神，成为我毕设时攻坚克难的重要支撑，让我在学术道路上更加坚定自信。

春晖寸草，难以回报。

我由衷感谢我的父母。是他们无言的坚守，伴我走过求学之路的风风雨雨。他们用无尽的爱与包容，为我遮风挡雨、默默奉献。他们几乎倾尽时间与精力，只为让我无忧无虑地专注于学业。无论我做出怎样的选择，他们始终给予我最无私的支持与鼓励。他们是我最坚实的后盾，也是我生命中最深的恩典。


山水一程，三生有幸

感谢我的室友李宇硕、陈涛、黄翊轩。四年同窗，我们共度朝暮，友情早已超越了“室友”二字。在我遭遇挫折时，是你们的鼓励和支持让我重新振作；在我疲惫迷茫时，是你们的陪伴让我找到归属。你们的笑语和真情，镌刻下了大学四年最温暖的篇章。愿我们前程似锦，不负芳华。

此外，我也要感谢机器人队的曲九合、付宇轩、华骏扬、彭会帅等同学。与你们共事一年多的时光里，我不仅学到了许多技术要领与实践经验，更在一次次合作与交流中，体会到团队的力量与同行的温暖。那些深夜调试的代码，激烈讨论的方案，并肩攻克的项目，都成为了我宝贵的回忆和成长的见证。希望未来我们都能继续在各自的道路上发光发热，不枉这段并肩奋战的时光。

感谢青葱岁月里与我同行的王英泰、刘廷汉、王锦翔、刘志博、左艺佳、洑淑琴等同学。无论是课余时光中畅谈理想、思维碰撞，还是课程设计时的协同创新、精益求精，你们的专业素养与执着精神始终让我受益匪浅。那些平日里撞出的灵感火花，不仅拓宽了我的视野，更让我感受到同行者的珍贵，这段携手共进的大学时光，将永远镌刻在我记忆深处。

深深感谢我的母校，北京理工大学。四年的大学时光中，学校为我提供了良好的学习环境与丰富的资源平台，让我有机会不断追求学术理想、提升综合素质。这里的老师和同学、教室与实验室，都将成为我人生中最宝贵的回忆。愿母校桃李芬芳满天下，英才辈出续华章。

最后，我要感谢那个始终不曾迷失方向的自己，那个不忘初心，熬过漫长昏暗时光的自己。纵然跌倒千百次，亦能倔强站起，不为荣耀，只因心中那束微光从未熄灭。这样的坚持与热望，带我走过风雨，也让我在岁月的渡口，得以无悔回望。君子坐而论道，少年当起而行，我愿初心不改，行动为先，持之以恒，终有所成。




\end{acknowledgements}
